%& -shell-escape
% !TeX root = thesis-abstract.tex
% !TeX encoding = UTF-8
% !TeX program = XeLaTeX

\documentclass{.local/lib/classes/ndvr-super-article-standalone}

\begin{document}

    %\pagebreak
    Метод поиска событий на основе 
    сравнения скользящих средних оценок сигнала
    выражается формулой:
    \[
        V \esim  W \Longleftarrow  \left\lbrace
            \begin{array}{rll}
                \abs{\overline{Y}_{A,t} - \overline{Y}_{B,t}} & \to & 0; \\
                \overline{Y}_{s,t} = \overline{Y}_{s,t,n}  & = & \frac{1}{s}\sum\limits^{t}_{k=t-s} Y_{k,n}.
            \end{array}
        \right.   
    \]
    Здесь:
    \begin{itemize}
        \item {
            $A$, $B$ — выбранные размеры скользящих окон;
        }
        \item {
            $Y_{k}$ — состояние видео в момент времени $k$.
        }
    \end{itemize}
    Для двумерного кадра с шириной $w$ пикселей и высотой $h$ 
    пикселей $Y_{k,n}$ задается как $
        Y_{k,n} = Y_{k,2}  =  \frac{1}{w \cdot h } 
            \sum\limits_{i=0}^{w} \sum\limits_{j=0}^{h}
                y_{k}^{(i, j)}
    $, где $y_{k}^{(i, j)}$ — яркость 
    пикселя $(i, j)$ в момент времени $k$.
    В общем случае, предлагается рассматривать $n$-мерный кадр.
    Таким образом, \[
            Y_{k,n} = 
            \dfrac
                { \sum\limits_{i_1=0}^{w_1} 
                    \ldots 
                        \sum\limits_{i_n=0}^{w_n} 
                            y_{k}^{(i_1\ \ldots\ i_n)}}
                {\prod\limits_{j=1}^{n} w_j  } .
    \]

    В процессе изучения методов выделения аномалий
    разработан декларативный алгоритмический 
    язык потоковой фильтрации.
    Фильтр сравнения скользящих средних оценок сигнала
    на~языке потоковой фильтрации приведён на~рисунке 
    \ref{fig:mean-sign-diff-event-filter}.

    Язык потоковой фильтрации применяется 
    как средство описания процесса обработки видео.
    Основная сущность языка — фильтр. 
    Фильтр — это абстрактный тип на языке Python\footnote{
        Guido van Rossum, Drake Fred L. 
        The Python Language Reference Manual. ––
        Bristol, United Kingdom: Network Theory Ltd., 
        2003. –– Sep. –– P. 144. 
        –– ISBN:1906966141, 9781906966140.
    } 
    с перегруженными операторами. 
    Фильтры описывают набор действий над~временным рядом.
    При любых операциях над~фильтрами вычисления 
    с потоком видео не производятся. 
    В этот момент происходит только создание дерева операций.
    Полное дерево операций выражено итоговым фильтром 
    («result\_filter»). 
    Непосредственная обработка видео происходит 
    в момент вызова метода 
    «filter\_objects» итогового фильтра «result\_filter».
    Метод «filter\_objects» обходит дерево в ширину
    и применяет операции из узлов дерева к потоку видео.
    Если операции не зависят друг от друга,
    то их результаты вычисляются параллельно. 

    \begin{figure}[h]
        \includegraphics[width=0.9\linewidth]
        {.local/vec/out/pdf/mean-sign-diff-event-filter-code.pdf}
        \caption{
             Фильтр сравнения скользящих средних оценок сигнала.\\
             Оператор «|» означает «конвейер». 
        }
        \label{fig:mean-sign-diff-event-filter}
    \end{figure}

    \pagebreak
\end{document}
 
